%%%%%%%%%%%%%%%%%%%%%%%%%%%%%%%%%%%%%%%%%%%%%%%%%%%%%%%%%%%%%%%%%%%%%%
% Lớp 12 | Các số đặc trưng ... ghép nhóm | Khoảng biến thiên, khoảng tứ phân vị của mẫu số liệu ghép nhóm
% Dạng 2D3?1-1: Công thức lý thuyết
% Số câu: 8  (sinh tự động bởi tach_cau.py; từ nay thầy sửa trực tiếp file này)
%%%%%%%%%%%%%%%%%%%%%%%%%%%%%%%%%%%%%%%%%%%%%%%%%%%%%%%%%%%%%%%%%%%%%%

%%%%%%%%%%%%%%%% TRẮC NGHIỆM 4 PHƯƠNG ÁN %%%%%%%%%%%%%%%%

% nguồn: Lop_12_Dai_so_Chuong_3
\begin{ex}%[2D3N1-1]%[ma:c838391b]
	Số đặc trưng nào không sử dụng thông tin của nhóm số liệu đầu tiên  và nhóm số liệu cuối cùng?
	\choice 
		{Khoảng biến thiên} 
		 {\True Khoảng tứ phân vị}
		{Phương  sai}
		{Độ lệch chuẩn} 
	 \loigiai{
	Số đặc trưng không sử dụng thông tin của nhóm số liệu đầu tiên và nhóm số liệu cuối cùng  là khoảng tứ phân vị.}
\end{ex}

% nguồn: Lop_12_Dai_so_Chuong_3
\begin{ex}%[2D3N1-1]%[ma:e55cbb83]
	Nếu thay tất cả các tần số trong mẫu số liệu ghép nhóm trên bằng 4 thì số đặc trưng nào  sau đây không thay đổi?
	 \choice
		{\True  Khoảng biến thiên}
		{Khoảng tứ phân vị} 
		{Phương  sai}
		{Độ lệch chuẩn} 
	 \loigiai{
	Nếu thay tất cả các tần số trong mẫu số liệu ghép nhóm trên bằng 4 thì số đặc  trưng không đổi là khoảng biến thiên.}
\end{ex}

% nguồn: Lop_12_Dai_so_Chuong_3
\begin{ex}%[2D3H1-1]%[ma:0a7cefcb]
	Các bạn học sinh lớp 10A  trả lời $40$ câu hỏi trong một bài kiểm tra. Kết quả được thống kê ở bảng sau:
	\begin{center}
		\begin{tabular}{|c|c|c|c|c|c|}
			 \hline
			Số câu trả lời đúng & {$[16; 21)$} & {$[21; 26)$} & {$[26; 31)$} & {$[31;  36)$} & {$[36; 41)$} \\
			 \hline
			Số học sinh & 4 &  11 & 12 & 9 & 4 \\
			 \hline
		\end{tabular}
	\end{center}
	Xác định  nhóm có tần số lớn nhất.
	\choice 
		{\True $[26;  31)$}
		{$[21;  26)$}
		 {$[31; 36)$}
		{$[36;  41)$}
	\loigiai{ 
		Nhóm có tần  số lớn nhất là $[26; 31)$.
	 }
\end{ex}


%%%%%%%%%%%%%%%% TRẢ LỜI NGẮN %%%%%%%%%%%%%%%%

% nguồn: Lop_12_Dai_so_Chuong_3
\begin{bt}%[2D3H1-1]%[ma:af16dbc2]
	Biết hàm số $F(x)=\dfrac{2 x^2-3 x+6}{x}$ là một nguyên hàm của hàm  số $$f(x)=\dfrac{a x^2+b x+c}{x^2} \quad \text{ (với $a$, $b$, $c$ là các số thực)}.$$
	Khi đó  giá trị $a+b-c$ bằng bao nhiêu?
	\shortans{8} 
	 \loigiai{
		Ta có $F'(x)=\left(2 x-3+\dfrac{6}{x}\right)'=2-\dfrac{6}{x^2}=\dfrac{2 x^2-6}{x^2}$  và $F'(x)=f(x)$.
		\\ 
		Suy ra $a=2$; $b=0$; $c=-6$. Do  đó $a+b-c=8$.
	 }
\end{bt}

% nguồn: Lop_12_Dai_so_Chuong_3
\begin{bt}%[2D3H1-1]%[ma:87c6da6f]
	Biết hàm  số $F(x)$ là một nguyên hàm của hàm số $f(x)=2^x+3 x^2$ và $F(1)=\dfrac{2}{\ln 2}$. Giá trị của $F(2)$ (làm tròn đến hàng phần mười) bằng
	 \shortans{12,8}
	\loigiai{ 
		Ta  có
		\[ 
		\displaystyle\int f(x)\mathrm{\,d}x 
		=\displaystyle\int\left(2^x+3  x^2\right)\mathrm{\,d}x
		=\dfrac{2^x}{\ln  2}+x^3+C.
		\] 
		Suy ra  $F(x)=\dfrac{2^x}{\ln 2}+x^3+C$.
		\\ 
		Ta có $F(1)=\dfrac{2}{\ln 2}+1+C=\dfrac{2}{\ln 2} \Rightarrow  C=-1$.
		 \\
		 Do đó $F(2)=\dfrac{2^2}{\ln 2}+2^3-1 \approx 12{,}8$.
	} 
\end{bt}

% nguồn: Lop_12_Dai_so_Chuong_3
\begin{bt}%[2D3T1-1]%[ma:ca9ca983]
	Một ô tô đang chạy với vận tốc $17{,}5$ m/s thì người lái xe đạp phanh. Từ thời điểm đó ô tô chuyển động với vận tốc $v(t)=\dfrac{35}{2}-\dfrac{7}{2} t$ (m/s), trong đó $t$ (tính bằng giây) là thời gian kể từ lúc bắt đầu đạp phanh. Quãng đường ô tô  di chuyển từ lúc đạp phanh đến khi dừng hẳn bằng bao nhiêu mét?
	 \shortans{43,8}
	\loigiai{ 
		Ô  tô dừng hẳn khi
		 \[
		v(t)=0  \Leftrightarrow \dfrac{35}{2}-\dfrac{7}{2} t=0
		\Leftrightarrow t=5  \quad \text{ (giây)}.
		\] 
		Quãng đường ô  tô di chuyển từ lúc đạp phanh là
		\[ 
		s(t)=\displaystyle\int v(t)\mathrm{\,d}t 
		 =\displaystyle\int\left(\dfrac{35}{2}-\dfrac{7}{2} t\right)\mathrm{\,d}t
		=\dfrac{35 t}{2}-\dfrac{7  t^2}{4}+C.
		\] 
		Ta có $s(0)=0$, suy ra  $C=0$. Khi đó $s(t)=\dfrac{35 t}{2}-\dfrac{7 t^2}{4}$.
		\\ 
		Do vậy $s(5)=\dfrac{35}{2} \cdot 5-\dfrac{7}{4}  \cdot 5^2=43{,}8$ m.
	 }
\end{bt}

% nguồn: Lop_12_Dai_so_Chuong_3
\begin{bt}%[2D3B1-1]%[ma:8a39cad2]
	Cho hàm  số $f(x) = x\sqrt{x}$ có một nguyên hàm là $F(x)$ thỏa $F(1) = 3$. Tính $F(4)$.
	 \shortans{15,4}
	 \loigiai{
		Ta có $F(x)=\displaystyle \int  x\sqrt{x}\mathrm{\,d}x=\displaystyle \int x^{\frac{3}{2}}\mathrm{\,d}x=\dfrac{2}{5}x^{\frac{5}{2}}+C$.\\
		Mà $F(1)=3\Rightarrow  \dfrac{2}{5}+C=3\Rightarrow C =\dfrac{13}{5}$.\\
		 Do đó $F(x)=\dfrac{2}{5}x^{\frac{5}{2}}+\dfrac{13}{5}$.\\
		Vậy  $F(4)=\dfrac{77}{5}=15{,}4$.
	 }
\end{bt}

% nguồn: Lop_12_Dai_so_Chuong_3
\begin{bt}%[2D3T1-1]%[ma:30dac1bf]
	Kí hiệu $h(x)$ là chiều cao của một cây (tính theo mét) sau khi trồng $x$ năm. Biết rằng sau năm đầu tiên cây cao $3$ m. Trong các năm tiếp theo, cây phát triển với tốc độ $h'(x)=\dfrac{\sqrt{2}}{x}$ (tính theo mét/năm). Chiều cao của cây đó sau $5$  năm (làm tròn đến hàng phần mười) bằng bao nhiêu mét?
	\shortans{5,3} 
	 \loigiai{
		Ta  có
		\[ 
		 \displaystyle\int h'(x)\mathrm{\,d}x
		 =\displaystyle\int \dfrac{\sqrt{2}}{x}\mathrm{\,d}x
		 =\sqrt{2} \ln |x|+C.
		\] 
		Suy ra  $h(x)=\sqrt{2} \ln |x|+C$.
		\\ 
		Ta có $h(1)=\sqrt{2}  \ln 1+C=3$ hay $C=3$.
		\\ 
		Do đó $h(5)=\sqrt{2} \ln 5+3 \approx  5{,}3(m)$.
	 }
\end{bt}
